%%%%%%%%%%  scan idx 388  |  folio 375  %%%%%%%%%%
\SGAPrintPage{375}%
\sgaStableTarget{sga6.e06.lem.1-10}\textbf{Lemma 1.10.} \uline{With the notation and hypotheses of} 1.1, \uline{let} \(\underline{F}\) and \(\underline{G}\) \uline{be two locally free sheaves of finite type on} \(X\), \uline{such that}:
\[
  \cl_{K^{\bullet}(X)}(\underline{F})=\cl_{K^{\bullet}(X)}(\underline{G});
\]
\uline{then there exists an integer} \(n_0\), \uline{such that for} \(n\geq n_0\), \uline{we have}:
\[
\begin{array}{ll}
\text{i)}& f_*(\underline{F}(n)),\ f_*(\underline{G}(n))\text{ \uline{are locally free of finite type on} }S;\\[1mm]
\text{ii)}& \cl_{K^{\bullet}(S)}\bigl(f_*(\underline{F}(n))\bigr)=\cl_{K^{\bullet}(S)}\bigl(f_*(\underline{G}(n))\bigr).
\end{array}
\]

Let \(A\) be the free abelian group generated by the classes (up to isomorphism) of locally free sheaves of finite type on \(X\). To say that, in \(K^{\bullet}(X)\), \(\underline{F}\) and \(\underline{G}\) have the same class is to say, by the construction of \(K^{\bullet}\), that there exist exact sequences of \(\underline{O}_X\)-Modules locally free of finite type:
\[
  0\longrightarrow \underline{L}_i\longrightarrow \underline{M}_i\longrightarrow \underline{N}_i\longrightarrow 0,
  \qquad i=1,\ldots,k,
\]
and integers \(n_i\), such that in \(A\) we have the equality:
\[
  (\underline{F})-(\underline{G})=\sum_i n_i\bigl((\underline{L}_i)-(\underline{M}_i)+(\underline{N}_i)\bigr).
\]
For every \(n\), tensoring with \(\underline{O}_X(n)\) defines an automorphism of \(A\), whence:
\[
  (\underline{F}(n))-(\underline{G}(n))=\sum_i n_i\bigl((\underline{L}_i(n))-(\underline{M}_i(n))+(\underline{N}_i(n))\bigr),
\]
the sequences \(0\to \underline{L}_i(n)\to \underline{M}_i(n)\to \underline{N}_i(n)\to0\) remaining exact.

Let \(B\) be the free abelian group generated by the classes (up to isomorphism) of quasi-coherent sheaves on \(S\), and \(A'\subset B\) the subgroup generated by the classes of locally free sheaves of finite type on \(S\). Direct image defines a homomorphism \(A\to B\), and we therefore still have:
%%%%%%%%%%  scan idx 389  |  folio 376  %%%%%%%%%%
\SGAPrintPage{376}%
\[
\begin{aligned}
  \bigl(f_*(\underline{F}(n))\bigr)-\bigl(f_*(\underline{G}(n))\bigr)
  ={}&\sum_i n_i\bigl((f_*(\underline{L}_i(n)))-(f_*(\underline{M}_i(n)))\\
      &\hspace{35mm}+(f_*(\underline{N}_i(n)))\bigr).
\end{aligned}
\]

By 1.3, for sufficiently large \(n\), the following conditions will hold:
\[
  \forall k\geq1,\ \forall i,
  \quad R^k f_*(\underline{L}_i(n))=0;
\]
\[
  \forall i,
  \quad f_*(\underline{L}_i(n)),\ f_*(\underline{M}_i(n)),\ f_*(\underline{N}_i(n)),\ f_*(\underline{F}(n)),\ f_*(\underline{G}(n))
  \text{ are locally free of finite type.}
\]

It follows that for sufficiently large \(n\), these elements all belong to \(A'\) and satisfy the relation written above, and that the sequences:
\[
  0\longrightarrow f_*(\underline{L}_i(n))\longrightarrow f_*(\underline{M}_i(n))\longrightarrow f_*(\underline{N}_i(n))\longrightarrow0
\]
are exact. This implies that in \(K^{\bullet}(S)\), \(f_*(\underline{F}(n))\) and \(f_*(\underline{G}(n))\) have the same class.

1.11. Suppose given an equation of linear dependence over \(K^{\bullet}(S)\) among \(1,\xi,\ldots,\xi^r\):
\sgaStableTarget{sga6.e06.eq.1-11-1}\begin{equation}\tag{1.11.1}
  f^*(x_0)+f^*(x_1)\xi+\ldots+f^*(x_r)\xi^r=0.
\end{equation}

We can write each of the \(x_j\) as the difference of the classes of two sheaves, namely: \(x_j=\cl(\underline{E}_j)-\cl(\underline{E}'_j)\), where \(\underline{E}_j\) and \(\underline{E}'_j\) are locally free of finite type on \(S\). Put
\[
  \underline{F}=\bigoplus_{j=0}^r f^*(\underline{E}_j)(j),
  \qquad
  \underline{F}'=\bigoplus_{j=0}^r f^*(\underline{E}'_j)(j).
\]

Equality (1.11.1) then reads:
\[
  \cl_{K^{\bullet}(X)}(\underline{F})=\cl_{K^{\bullet}(X)}(\underline{F}').
\]

By Lemma 1.10, we obtain for sufficiently large \(n\):
\[
  \cl_{K^{\bullet}(S)}\bigl(f_*(\underline{F}(n))\bigr)=\cl_{K^{\bullet}(S)}\bigl(f_*(\underline{F}'(n))\bigr).
\]

%%%%%%%%%%  scan idx 390  |  folio 377  %%%%%%%%%%
\SGAPrintPage{377}%
Now
\[
\begin{aligned}
  f_*(\underline{F}(n))
  &=\bigoplus f_*\bigl(f^*(\underline{E}_j)\otimes_{\underline{O}_X}\underline{O}_X(n+j)\bigr)\\
  &=\bigoplus \underline{E}_j\otimes_{\underline{O}_S} f_*(\underline{O}_X(n+j)).
\end{aligned}
\]

Moreover, \(f_*(\underline{O}_X(n+j))=\Sym(\underline{E})^{n+j}\); we shall denote by \(\sigma_p\) the class in \(K^{\bullet}(S)\) of \(\Sym(\underline{E})^p\). Relation (1.11.1) therefore finally gives:
\sgaStableTarget{sga6.e06.eq.1-11-2}\begin{equation}\tag{1.11.2}
  \forall n\geq n_0,
  \qquad
  \sum_{j=0}^r x_j\sigma_{n+j}=0.
\end{equation}

We then consider the complex:
\[
  0\longrightarrow \bigwedge^{r+1}(\underline{E})\otimes_{\underline{O}_S}\Sym(\underline{E})
  \longrightarrow \bigwedge^r(\underline{E})\otimes_{\underline{O}_S}\Sym(\underline{E})
  \longrightarrow\ldots\longrightarrow
  \underline{E}\otimes_{\underline{O}_S}\Sym(\underline{E})
  \longrightarrow \Sym(\underline{E})\longrightarrow0,
\]
whose arrows are defined by:
\[
\begin{aligned}
&(x_1\wedge\ldots\wedge x_n)\otimes(y_1\otimes\ldots\otimes y_p)\rightsquigarrow{}\\
&\qquad \sum(-1)^{i-1}(x_1\wedge\ldots\wedge\widehat{x_i}\wedge\ldots\wedge x_n)
\otimes(x_i\otimes y_1\otimes\ldots\otimes y_p)
\end{aligned}
\]
If we regard \(\bigwedge^n(\underline{E})\) as a graded \(\underline{O}_S\)-Module concentrated in degree \(n\), this is a complex of graded \(\Sym(\underline{E})\)-Modules, the morphisms themselves being graded.

This complex is a resolution of \(\underline{O}_S\): on an open set where \(\underline{E}\) is free, \(\Sym(\underline{E})\) identifies with a polynomial Algebra \(\underline{O}_S[T_0,\ldots,T_r]\), and the complex with the Koszul complex defined by the \(\Sym(\underline{E})\)-regular sequence \(T_0,\ldots,T_r\); we know that the latter is a resolution of \(\Sym(\underline{E})/(T_0,\ldots,T_r)=\underline{O}_S\).

Writing that each homogeneous component of the complex of degree \(\geq1\) is acyclic, we therefore obtain relations in \(K^{\bullet}(S)\):
\sgaStableTarget{sga6.e06.eq.1-11-3}\begin{equation}\tag{1.11.3}
  \sum_{i=0}^{r+1}(-1)^i\lambda^i(E)\sigma_{p-i}=0
  \qquad\text{for every }p\geq1.
\end{equation}
(putting \(\sigma_k=0\) for \(k<0\)).

Suppose relation (1.11.2) holds for \(n\geq n_0\); we shall show that it still holds for \(n_0-1\), which will allow us to descend by induction to \(n_0=-r\) (we shall see later that this suffices to imply the vanishing of the \(x_j\)).

%%%%%%%%%%  scan idx 391  |  folio 378  %%%%%%%%%%
\SGAPrintPage{378}%
By (1.11.3), we can write
\begin{equation}\tag{1.11.4}
  \sum_{j=0}^r x_j\left(\sum_{i=0}^{r+1}(-1)^i\lambda^i(E)\sigma_{p(j)-i}\right)=0,
\end{equation}
with \(p(j)\geq1\). We take \(p(j)=n_0+r+j\), which is indeed \(\geq1\) if \(n_0>-r\). We then obtain:
\[
\begin{aligned}
\sum_{j=0}^r x_j\left(\sum_{i=0}^{r+1}(-1)^i\lambda^i(E)\sigma_{n_0+r+j-i}\right)
&=\sum_{i=0}^{r+1}(-1)^i\lambda^i(E)
   \left(\sum_{j=0}^r x_j\sigma_{n_0+r+j-i}\right).
\end{aligned}
\]

For \(i\neq r+1\), \(n_0+r-i\geq n_0\), so by (1.11.2):
\[
  \sum_{j=0}^r x_j\sigma_{n_0+r+j-i}=0.
\]

Consequently, relation (1.11.4) reduces to:
\[
  \lambda^{r+1}(E).\left(\sum_{j=0}^r x_j\sigma_{n_0-1+j}\right)=0.
\]

Now \(\underline{E}\) has rank \(r+1\), so \(\lambda^{r+1}(E)\) is the class of an invertible sheaf, and is therefore invertible in \(K^{\bullet}(S)\). Consequently, relation (1.11.4) reduces to relation (1.11.2) for \(n_0-1\).

Thus relation (1.11.2) holds for \(n\geq-r\); for \(n=-r\), it gives:
\[
  x_r\sigma_0=0.
\]
Since \(\sigma_0=1\), we obtain \(x_r=0\). We thus show by descending induction on \(j\) that the \(x_j\) are all zero, using relations (1.11.2) for \(n\) ranging from \(-r\) to \(0\). This shows that \(1,\ldots,\xi^r\) are independent over \(K^{\bullet}(S)\), and completes the proof of Theorem 1.1.

\sgaStableTarget{sga6.e06.rem.1-12}\textbf{Remark 1.12.} For every \(\underline{O}_S\)-Module locally free of finite type \(\underline{E}\), let \(\sigma_i(E)\) be the class in \(K^{\bullet}(S)\) of the \(i\)-th symmetric power of \(\underline{E}\); put \(\sigma_t(E)=1+\sum_{i\geq1}\sigma_i(E)t^i\); relations (1.11.3) can also be written:
%%%%%%%%%%  scan idx 392  |  folio 379  %%%%%%%%%%
\SGAPrintPage{379}%
\begin{equation}\tag{1.12.1}
  \lambda_t(E).\sigma_{-t}(E)=1.
\end{equation}

This relation itself shows the additivity of the map \(\sigma_t\), which can therefore be extended to a homomorphism from \(K^{\bullet}(S)\) to \(1+K^{\bullet}(S)[[t]]^+\):

\sgaStableTarget{sga6.e06.rem.1-13}\textbf{Remark 1.13.} We can take as a basis of \(K^{\bullet}(X)\) over \(K^{\bullet}(S)\) any sequence of powers of \(\xi\); \(\xi^k,\ldots,\xi^{k+r}\), since \(\xi\) is invertible, and such a sequence therefore consists of elements linearly independent over \(K^{\bullet}(S)\); moreover, the other powers of \(\xi\) can be expressed explicitly as linear combinations of these with coefficients in \(K^{\bullet}(S)\), by equality (1.9.1) and the invertibility of \(\xi\).

Put \(x=1-\xi^{-1}\) (\(x\) is interpreted as the ``class of a hyperplane section''). We can then take as a basis of \(K^{\bullet}(X)\) over \(K^{\bullet}(S)\) the elements \(1,x,\ldots,x^r\). To obtain the equation of dependence of \(1,\ldots,x^{r+1}\) we consider the exact sequence:
\[
  0\longrightarrow \underline{F}\longrightarrow f^*(\underline{E})\longrightarrow \underline{O}_X(1)\longrightarrow0,
\]
and the exact sequence obtained by passing to the duals:
\[
  0\longrightarrow \underline{O}_X(-1)\longrightarrow f^*(\check{\underline{E}})\longrightarrow \check{\underline{F}}\longrightarrow0.
\]
We then observe that \(\gamma^{r+1}(\check F-r)=0\), from which we deduce:
\[
  \gamma^{r+1}(\check F-r)=\gamma^{r+1}(\check E-r-1+x)=0,
\]
that is:
\sgaStableTarget{sga6.e06.eq.1-13-1}\begin{equation}\tag{1.13.1}
  \sum_{k=0}^{r+1}\gamma^{r+1-k}(\check E-r-1).x^k=0.
\end{equation}

We obtain an analogous result by introducing \(x'=\xi-1\).

Finally, note that if \(\underline{E}=(\underline{O}_S)^{r+1}\), and hence \(X=\mathbf P_S^r\), we have \(E=r+1\), and equality (1.13.1) reduces to:
